\section{本章小结}
\label{sec:integral-summary}

本章回答的问题是：给定测度，怎样累计一个函数的取值，以及累计之后哪些运算可以交换。积分从简单函数的分块累计开始（第~\ref{sec:integral-definition}节）：在取值固定的可测块上，贡献就是取值乘以块的测度。非负可测函数用下方简单函数的积分上确界定义，二进制分层逼近说明这个上确界确实能逐层读出函数；可正可负的函数拆成正部和负部，复值函数拆成实部和虚部，绝对可积条件保证不出现$\infty-\infty$；对有符号测度则沿Jordan分解分别积分。整个构造只用到测度和被积函数本身，不需要先寻找哪个测度的密度。随后，单调收敛定理把简单函数的运算规则传递到一般函数（第~\ref{sec:integral-properties}节），离散求和、闭区间上的Riemann积分和绝对收敛的反常积分都成为同一积分的特例（第~\ref{sec:integral-classical}节）。

本章的大部分结论都在回答“什么条件允许交换两种运算”，而起作用的条件可以按运算归纳为表~\ref{tab:integral-summary-conditions}。非负性和单调递增支持递增累计，结果允许为无穷；共同的可积控制防止质量逃逸，支持极限与积分交换（第~\ref{sec:integral-limits}节）；绝对可积性防止正负贡献以不同方式抵消，支持改变累次积分的次序（第~\ref{sec:integral-iterated}节）。换元则按映射的条件逐层具体化：一般可测映射只需推前测度，可逆光滑变换由Jacobian行列式给出体积比例，分支可逆变换要累计所有原像分支（第~\ref{sec:integral-change-of-variables}节）。含参积分把参数变化转成函数列的极限，连续性与积分号下求导都要求在参数点附近有统一的可积控制（第~\ref{sec:integral-parametric}节）。

\begin{table}[!htbp]
\centering
\caption{本章主要运算及其成立条件。非负情形的等式允许两侧同为$+\infty$；表中条件都是充分条件，反例说明去掉条件后结论可能失败。}
\label{tab:integral-summary-conditions}
\small
\begin{tabularx}{\linewidth}{@{}l>{\RaggedRight\arraybackslash}X>{\RaggedRight\arraybackslash}X@{}}
\toprule
运算 & 条件 & 依据与反例\\
\midrule
极限与积分交换 & 非负递增 & 单调收敛\\
 & 几乎处处收敛，有可积控制 & 控制收敛；$n\mathbf 1_{(0,1/n)}$\\
累次积分换序 & 非负，两侧$\sigma$有限 & Tonelli；对角线反例\\
 & 乘积测度下绝对可积 & Fubini；有理函数反例\\
换元 & 可测映射 & 推前换元\\
 & 微分同胚或可数分支 & Jacobian公式；分支累计\\
参数连续、求导 & 参数点附近统一控制 & 控制收敛；$te^{-tx}$\\
由导数恢复函数 & 绝对连续 & 基本定理；Cantor函数\\
\bottomrule
\end{tabularx}
\end{table}

积分与微分的关系在第~\ref{sec:integral-differentiation}节收束。可积函数的不定积分绝对连续，并且在几乎每一点的导数等于原函数；反过来，绝对连续函数恰好是其导数的不定积分加常数。两个方向分别依靠积分在小测度集上的绝对连续性、单调函数的几乎处处可导，以及区间覆盖下的零导数引理。导数可积与能够由导数恢复函数是两回事：阶跃函数和Cantor函数都几乎处处可导、导数可积，却不能由导数恢复；后者连续单调，增量全部落在零测的Cantor集上。

下一章把问题反过来。本章固定一个测度，用函数的积分产生数值；如果给定两个测度，能否找到一个函数，使它相对于参考测度的积分恢复目标测度在每个集合上的质量？Cantor函数已经显示这样的函数未必存在。第四章将以本章的积分定义、单调收敛定理和基本运算为依据，研究这种密度何时存在、怎样计算，以及如何用它在同一空间上转换积分。
